\scanpage{50}

\uline{Proposition} Soit $\Pi$ un tétraèdre orienté, $R$ $(R^{+})$ l'ens de ses repères (\inserted{resp.} directs), $S,A,F$ l'ens des ses sommets, \deleted{\unsure{a}}arêtes, faces -- d'où $\rho_s,\rho_a,\rho_f$ si $s\in S$, $a\in A$, $f\in F$. \deleted{Soit} $\Phi=\underset{\inserted{4}}S\amalg\underset{\inserted{6}}A\amalg\underset{\inserted{4}}F$ l'ens des facettes

\deleted{a) Les applic$^{\mathrm{ns}}$ $s\unsure{\to|\mathord{-}}\rho_s$,}

a) L'application $\varphi\unsure{\longrightarrow|\mathord{-}|\mapsto}\rho_{\unsure{p|\varphi}}$ de $S\amalg F$ dans $H=\mathrm{Aut}^{+}\Pi$ est injective, et son image est l'ens des $8$ él. d'ordre $3$ de $H$. De plus, si $s\in S$, $f\in F$ sont opposés, on a \inserted{$\rho_s\rho_f=\rho_f\rho_s=1$}

b) L'application $a\longmapsto\rho_a$ de $A$ dans $H$ induit une application de degré $2$ de $A$ dans \inserted{(sur)} l'ens $ {}_3H$ des él. d'ordre $3$ de $H$ [\uline{NB} deux arêtes sont dites opposées si elles ont $=$ image, \deleted{\unsure{en}} qui signifie aussi que ce sont des parties disjointes de $S$, ou de $F$\unsure{]|.}

c) L'application $r\unsure{\longrightarrow|\mathord{-}|\mapsto}\Sigma_r$ de $R$ dans l'ens des syst. de Schwartz de $H$ est injective [en effet le commutant de $H$ dans $\inserted{\widetilde H=}\mathrm{Aut}(\Pi)$ est $1$] l'image forme donc une classe sous $\widetilde H$ (\deleted{qui} c'est un torseur sous $\widetilde H$) et deux classes sous $H$, de type numérique $(3,3,2)$.

\margin{$24=12+12$ syst. de Schwartz}

% slanted
\margin{[celà fait $72=36+36$ en transformant par $\unsure{\mathfrak S_3|\mathcal S_3}$, deux paquets de $36$ syst. de Schwartz, \deleted{\unsure{chacun des} paquets de $72$ est stable par $\mathfrak S_3$ \unsure{et} \illegible{1}} stables sous $\mathfrak S_3\times H$}

d) Si on a un syst. de Schwartz canonique $(\rho_f,\unsure{\rho'_s|\rho_s},\rho'_a)$ \inserted{[ou anticanonique $(\rho'_{f},\rho'_{s},\rho'_{a})$]}, les $6$ systèmes qu'il engendre, par \deleted{\unsure{son|action}} \deleted{\unsure{$\rho_f,\rho_s,\rho_a$}} \deleted{$(\rho,\rho',\rho'')\ (\rho'',\rho,\rho')$} \deleted{\unsure{cla}} gpent donc $2$ par $2$ suivant des \deleted{$(\rho,\rho',\rho''),(\rho',\rho'',\rho)(\rho'',\rho,\rho')($} types numériques égaux, ceux qui sont de type numérique comme

\[
(\rho_f,\rho_s,\rho_a),\quad
(\underset{\inserted{\rho''_{f'}}}{\rho_s^{-1}},
 \underset{\inserted{\rho''_{s'}}}{\rho_f^{-1}},
 \underset{\inserted{\rho''_{a'}}}{\rho_a^{-1}=\rho_a})
\qquad\text{i.e. }\Sigma_r\text{ et }\Sigma_{r'}
\]

($r'$ le repère \uline{opposé} de $r$, il \deleted{\unsure{a}} a \uuline{$=$} orientation !) \deleted{sont conj} sont \deleted{\unsure{non}} conjugués par $H$, \deleted{mais conjugués par} \inserted{à fortiori par} $\widetilde H$. \deleted{(car $r,r'$ définissent des orientations opposées), donc on trouve trois classes de conjugaison distinctes, chacune}

% archive sheet number: 49
